% fig76-borders.tex -- Figure 14.x: borders of s = abcababcab (n = 10): longest border abcab (5); its border ab (2).
% fstyle.tex -- shared TikZ styles for the figures of this book (palette from lifeos.sty).
\tikzset{%
  cell/.style={draw=ink, minimum width=7mm, minimum height=6.5mm, inner sep=0pt, font=\footnotesize, fill=white},
  cellh/.style={cell, fill=accentsoft, draw=accent},
  cella/.style={cell, fill=amberbg, draw=amber},
  cellg/.style={cell, fill=goodbg, draw=good},
  cellx/.style={cell, fill=graybg, draw=slate, text=slate},
  vtx/.style={circle, draw=ink, minimum size=6.5mm, inner sep=0pt, font=\footnotesize, fill=white},
  vtxh/.style={vtx, draw=accent, fill=accentsoft},
  vtxa/.style={vtx, draw=amber, fill=amberbg},
  vtxg/.style={vtx, draw=good, fill=goodbg},
  vtxd/.style={vtx, draw=slate, fill=graybg, text=slate},
  lab/.style={font=\footnotesize, text=slate},
  labA/.style={font=\footnotesize, text=accent},
  labM/.style={font=\footnotesize, text=amber},
  idx/.style={font=\tiny, text=slate},
  arr/.style={-{stealth}, thick, ink},
  arrA/.style={-{stealth}, thick, accent},
  arrM/.style={-{stealth}, thick, amber},
  arrG/.style={-{stealth}, thick, good},
  edge/.style={thick, ink},
  edgeA/.style={thick, accent},
  edgeM/.style={thick, amber},
}%
% \cellrow{prefix}{x0}{y}{style}{v0,v1,...}: a row of cells named prefix0, prefix1, ... (pitch 0.75cm)
\providecommand{\cellrow}[5]{\foreach \v [count=\i from 0] in {#5} {\node[#4] (#1\i) at (#2+\i*0.75,#3) {\v};}}%
% \cellidx{prefix}{n}{start}: index labels under cells prefix0 .. prefix(n-1)
\providecommand{\cellidx}[3]{\foreach \i in {0,...,\numexpr#2-1\relax} {\node[idx, below=0pt of #1\i] {\the\numexpr\i+#3\relax};}}%
% \fitfig{<tikzpicture>}: shrink the picture to \linewidth only when it would overflow.
\ifdefined\fitbox\else\newsavebox{\fitbox}\fi%
\providecommand{\fitfig}[1]{\sbox{\fitbox}{#1}\ifdim\wd\fitbox>\linewidth\resizebox{\linewidth}{!}{\usebox{\fitbox}}\else\usebox{\fitbox}\fi}%

\begin{tikzpicture}[font=\footnotesize]
\foreach \c [count=\i from 0] in {a,b,c,a,b,a,b,c,a,b}
  \node[\ifnum\i<5 cellh\else cella\fi, minimum width=7mm] (s\i) at (\i*0.78,0) {\c};
\draw[decorate,decoration={brace,raise=3pt},accent,thick] (s0.north west) -- (s4.north east) node[midway,above=5pt,labA]{prefix \texttt{abcab}};
\draw[decorate,decoration={brace,mirror,raise=3pt},amber,thick] (s5.south west) -- (s9.south east) node[midway,below=5pt,labM]{suffix \texttt{abcab}: border of length $5$};
\draw[good, very thick] (s0.north west) rectangle (s1.south east); \draw[good, very thick] (s8.north west) rectangle (s9.south east);
\node[font=\footnotesize, text=good, anchor=north] at (s9.south |- 0,-1.55) {the inner border \texttt{ab} (length $2$) is outlined in green};
\node[lab, anchor=west] at (0,-2.4) {chain of borders: $\pi_{10}=5$, then the border of \texttt{abcab} is $\pi_5=2$, then $\pi_2=0$};
\end{tikzpicture}
